% =============================================================================
\section{Introduction}\label{sec:introduction}
% =============================================================================
\IEEEPARstart{M}{agnetic} resonance imaging is intrinsically complex-valued. Receiver coils measure oscillating electromagnetic signals whose samples have real and imaginary components. Conventional image formation eventually produces magnitude images that are convenient for visualization and clinical interpretation, but the magnitude operation is many-to-one: different complex signals can map to the same nonnegative intensity. A diagnostic network trained only on magnitude therefore never receives the phase-bearing part of the acquisition, even when that information is present in the stored multicoil signal.

A complex sample $z=a+\mathrm{i}b=r\mathrm{e}^{\mathrm{i}\phi}$ jointly represents amplitude $r$ and phase $\phi$. Complex-valued processing is not merely a two-channel implementation choice. Complex multiplication constrains how the Cartesian components interact, global rotations have a natural group action, and several nonlinearities and pooling rules can be distinguished by whether they preserve, select, or explicitly modify phase. A \ac{CVNN} therefore imposes a structured inductive bias that differs from both magnitude-only learning and unconstrained real-valued processing of real and imaginary channels.

Complex-valued deep learning has been developed extensively at the methodological level. Trabelsi \emph{et al.} introduced practical complex convolution, normalization, initialization, and activation machinery for deep networks~\cite{Trabelsi2018DeepComplex}. In MRI, complex-valued models have been studied mainly in reconstruction and phase-sensitive applications. Cole \emph{et al.} performed parameter-controlled analyses for MRI reconstruction and phase-focused tasks~\cite{Cole2021ComplexMRI}; Xiao \emph{et al.} used complex-valued networks for partial-Fourier recovery of complex MR images~\cite{Xiao2022ComplexPartialFourier}; and Virtue \emph{et al.} examined phase-sensitive complex activations for MR fingerprinting~\cite{Virtue2017ComplexActivation}. These studies establish that complex representation can be useful when the target is close to the acquisition physics. Diagnostic classification is different: the label is several transformations removed from the measured signal, and the phase-bearing information may be useful, irrelevant, or confounded by scanner and coil effects.

The fastMRI Prostate release creates an unusually suitable setting for studying this question because it combines raw multicoil MRI data, reconstructed images, and slice-, volume-, and exam-level annotations in a public prostate MRI cohort~\cite{Tibrewala2024FastMRIProstate}. The release contains axial T2-weighted and diffusion-weighted acquisitions from clinical 3~T systems and provides research PI-RADS labels. The present work deliberately focuses on the prepared T2-weighted branch. This restriction is not a claim that T2-weighted imaging is sufficient for clinical prostate assessment. PI-RADS uses complementary sequence information and sequence dominance depends on anatomical zone~\cite{Turkbey2019PIRADSv21}. Rather, T2-weighted data provide a controlled substrate for the present signal-representation question: they preserve multicoil complex structure, depict prostate zonal anatomy and intraglandular morphology, and avoid additional diffusion-specific transformations such as EPI distortion handling, multiple $b$-values, directional averaging, ADC calculation, and derived high-$b$-value images~\cite{Tibrewala2024FastMRIProstate}.

The classification target is also deliberately narrow. Each sample is one axial T2-weighted slice and the released slice-level PI-RADS annotation is grouped into low-suspicion PI-RADS 1--2 versus higher-suspicion PI-RADS 3--5. This follows the threshold used in the dataset's technical classification demonstration~\cite{Tibrewala2024FastMRIProstate}, but the task should not be read as histopathological cancer diagnosis. PI-RADS 3 is an intermediate category, and the released slice labels were assigned retrospectively in isolation rather than from the complete multiparametric examination. The task is therefore best viewed as slice-level discrimination of low versus elevated radiological suspicion under a fixed research labeling protocol.

\subsection{Why a Factorial Complex-Valued Study Is Needed}
A comparison between one magnitude CNN and one arbitrarily chosen CVNN cannot identify why a difference occurs. Complex architectures involve several mathematically consequential choices. Standard complex convolution imposes complex linearity, whereas widely-linear convolution introduces a conjugate pathway and spans a broader real-linear class. RMS-type normalization preserves the orientation of a complex activation while scaling its power, whereas complex BatchNorm whitens the real-imaginary covariance. Magnitude-gated nonlinearities preserve phase direction, Cartesian ReLU does not, and cardioid gating explicitly depends on phase. Pooling can select an observed complex value by magnitude or combine values through vector averaging. Finally, a network may use only a complex stream or may combine complex features with a magnitude stream through late fusion or explicit cross-stream interactions.

These choices are not cosmetic hyperparameters. They encode different algebraic assumptions about how useful information should transform under global phase rotation, how much real-linear freedom the model should have, and how local phase relationships should survive downsampling. The present study therefore treats the architecture itself as an experimental object. We construct an exhaustive $2\times3\times2\times2\times4\times5$ grid of 480 complex models and match every complex configuration to the scalar parameter budget of a fixed real magnitude baseline.

\subsection{Scope and Contributions}
The study asks a controlled question: given the same patient-disjoint data partitions and approximately the same scalar parameter budget, how does performance vary when the network is allowed to retain complex information under different architectural constraints? The paper does not attempt to prove that MRI phase is a biological biomarker, nor does it evaluate a full clinical biparametric diagnostic pipeline. Its contribution is a structured experimental characterization of complex-valued representation and inductive bias for a fixed T2-weighted slice-classification task.

The main contributions are:
\begin{itemize}
    \item We formulate the magnitude map, global phase action, Fourier-domain representation, standard complex linearity, and widely-linear relaxation in a common algebraic framework for diagnostic MRI classification.
    \item We give an analytic characterization of the factorial operators---normalization, activation, pooling, stream fusion, and holographic interaction---in terms of real-linear freedom, phase dependence, and global-rotation behavior.
    \item We define a patient-disjoint slice-level T2-weighted PI-RADS-group classification task from the fastMRI Prostate data and explain the rationale and limitations of using T2-weighted slices as the controlled input sequence.
    \item We perform a parameter-matched factorial study containing one real magnitude baseline and 480 complex configurations, with all numerical summaries regenerated directly from the completed experiment CSV.
    \item We report both discrimination and minority-class metrics and explicitly separate descriptive factorial characterization from claims that would require repeated training realizations or independent external validation.
\end{itemize}

The remainder of the paper is organized as follows. Section~\ref{sec:math} develops the algebraic and analytic formulation. Section~\ref{sec:data} describes the dataset, T2-weighted scope, labels, preprocessing, and classification task. Section~\ref{sec:models} defines the model families and factorial operators. Section~\ref{sec:design} specifies the parameter-matched experiment and evaluation protocol. Section~\ref{sec:results} reports the observed factorial surface, followed by discussion and conclusions in Sections~\ref{sec:discussion} and~\ref{sec:conclusion}. The appendices collect material that can later be moved to a supplement or expanded arXiv version.
